\documentclass[11pt,a4paper]{article}
\usepackage[margin=24mm]{geometry}
\usepackage{amsmath,amssymb,amsthm,booktabs,microtype}
\setlength{\emergencystretch}{2em}
\newtheorem{theorem}{Theorem}
\newtheorem{lemma}[theorem]{Lemma}
\newcommand{\GL}{\operatorname{GL}}
\newcommand{\SL}{\operatorname{SL}}
\newcommand{\PGL}{\operatorname{PGL}}
\newcommand{\PSL}{\operatorname{PSL}}
\title{The remaining quotients and the linear polytope classifications}
\author{Human Compression project: completion of the N012 arguments}
\date{Component pending distinct review}
\begin{document}\maketitle
We complete the small quotient cases in Markovskaya's thesis [T] and
deduce its finite linear polytope classification from the constructions
already recovered. All dimensions below satisfy $n\ge2$. A group is
called admissible when three involutions generate it and two of those
involutions commute, with the source convention permitting repetitions.
This is weaker than having a rank-three string C-marking.

For odd $q$ put
\[
 H_n(q)=\{g\in\GL_n(q):\det g\in\{1,-1\}\},\qquad
 Q_n(q)=H_n(q)/Z(\SL_n(q)).
\]
The quotient here is the source's $\PGL_n^{\pm1}(q)$; dividing by
the full centre of $H_n(q)$ is in general a different operation.

\section{Precisely retained theory}
We acquire the complete proofs in four companion components:
\emph{Integral linear groups generated by three involutions} [I],
\emph{Three involutions in finite general linear groups} [L],
\emph{Central quotients of the determinant-sign subgroup} [Z], and
\emph{Involutions and intersections in groups with determinant
$\pm1$} [H]. The common root and fixed-space proofs accompanying [L]
are acquired as well. The following external inputs remain explicit:
\begin{enumerate}
\item The admissibility classification for simple projective linear
groups retained in [L]. For odd $q$, $\PSL_2(q)$, $q\ge5$, is
admissible except when $q=7,9$; every $\PSL_3(q)$ is excluded;
$\PSL_n(q)$ is admissible for $n\ge4$. In even characteristic,
the exclusions in dimensions three and four and the positive cases
in dimensions at least five are retained too.
\item The Sjerve--Cherkassoff theorem that $\PGL_2(q)$ is admissible
for every $q>2$, already retained in [L].
\item The Conder--Oliveros criterion and the standard correspondence
between string C-groups and regular polytopes, as retained in [L,Z].
For a finite group with no nontrivial cyclic normal subgroup, three
distinct generating involutions with a commuting pair satisfy the
rank-three intersection property.
\end{enumerate}
These are borrowed results, not new proofs of the external
classifications. Their full external proof and software costs remain
unmeasured. The integral existence theorem retained externally by [L]
is supplied here by the full proof [I]. There is no circular use:
[I] uses only the separately stated finite exclusions
$\PSL_3(2),\PSL_4(2)$, not the finite general-linear classification
that subsequently uses [I].

\section{Two and three coordinates}
The centre calculation proved in [Z] says that, with
$d=\gcd(n,q-1)$ and $e=(q-1)/d$,
\[
 Q_n(q)\cong
 \begin{cases}
 \PSL_n(q)\times C_2,&e\text{ even},\\
 \PSL_n(q)\rtimes C_2\le\PGL_n(q),&e\text{ odd}.
 \end{cases}                                            \tag{1}
\]
In the second case the embedding is the actual projective map.
We also use [Z]'s complete central-lift lemma: a rank-three string
C-marking of a perfect finite group $G$ lifts to one of $G\times C_2$
by giving all three generators the nonidentity second coordinate.
Its proof establishes both generation and the intersection property.

\begin{lemma}
For odd $q$, the group $Q_2(q)$ is admissible if and only if $q\ne9$.
In precisely those cases it is the automorphism group of a regular
polytope of rank three. No $Q_3(q)$, for odd $q$, is admissible.
\end{lemma}
\begin{proof}
When $n=2$, we have $d=2$. If $q\equiv3\pmod4$, then $e$ is odd.
The projective image in (1) contains $\PSL_2(q)$ with index two.
Since $\PGL_2(q)$ itself has index two over $\PSL_2(q)$, the image
is all of $\PGL_2(q)$. This is also seen directly from determinants:
$-1$ is a nonsquare, so the determinant-minus-one coset maps to the
nontrivial projective determinant coset.

For $q=3$, the action on the four projective points identifies this
group with $S_4$. The adjacent transpositions
$(12),(23),(34)$ give a string C-marking: the first and third commute,
and the subgroups $S_{\{1,2,3\}}$ and $S_{\{2,3,4\}}$ intersect
in $\langle(23)\rangle$. This covers the one nonsimple socle here.

For $q\equiv3\pmod4$, $q\ge7$, the retained theorem supplies
admissibility of $\PGL_2(q)$. Its nonabelian simple socle is
$S=\PSL_2(q)$. A cyclic normal subgroup $N$ intersects $S$ trivially:
the intersection is normal in $S$, while $S$ is not cyclic. Therefore
$[N,S]\le N\cap S=1$. The projective centralizer of $S$ is trivial.
Indeed, a representative of a projective centralizer conjugates
every transvection to a scalar multiple of itself. Both conjugate
matrices have sole eigenvalue one, so that scalar is one. Commuting
with the two off-diagonal matrix units then forces the representative
to be scalar. Hence $N=1$. A degenerate admissible triple would leave
at most two involutions and generate a dihedral group, which cannot
contain the nonabelian simple group $S$. Thus the three involutions
are distinct and nonidentity. The retained intersection criterion
gives a rank-three string C-marking.

If $q\equiv1\pmod4$, (1) gives
$Q_2(q)\cong\PSL_2(q)\times C_2$. Here $q\ge5$ and the simple
group is perfect. Its only inadmissible case in this congruence class
is $q=9$, by the retained classification; $q=7$ is in the other class.
For $q\ne9$, that classification and the intersection criterion
give a string C-marking of $\PSL_2(q)$. The central-lift lemma gives
one for $Q_2(q)$. At $q=9$, projection onto the inadmissible factor
$\PSL_2(9)$ rules out admissibility of the product, since admissibility
passes to quotients.

For $n=3$, $d$ is odd, so $e$ is even. Equivalently, the scalar $-I$
has determinant $-1$. Thus $Q_3(q)\cong\PSL_3(q)\times C_2$.
Projection onto $\PSL_3(q)$ and its retained nonadmissibility give
the final exclusion.
\end{proof}

\begin{theorem}
For odd $q$ and $n\ge2$, each of the following conditions is equivalent
to $n=2,q\ne9$, or $n\ge4$:
\begin{enumerate}
\item $Q_n(q)$ is admissible;
\item $Q_n(q)$ is the automorphism group of a regular polytope of rank three.
\end{enumerate}
\end{theorem}
\begin{proof}
The lemma settles dimensions two and three. The complete
four-dimensional theorem of [Z] gives the positive result for every
odd field, using its corrected matrices and its separate treatments
of the two centre cases. For $n\ge5$, [H] supplies an actual string
C-marking of every scalar quotient of $H_n(q)$; its kernel here is
the scalar subgroup $Z(\SL_n(q))$. All the claimed positive cases
therefore have regular polytopes. Conversely a rank-three string
C-marking is in particular an admissible triple, so the negative
cases exclude both conditions. This proves [T, Theorem 4.3.1 and
Corollary 4.3.1], with the source's actual quotient definition.
\end{proof}

\section{The finite general-linear polytope result}
\begin{theorem}
For $n\ge2$, the group $\GL_n(q)$ is the automorphism group of a
regular polytope of rank three if and only if $q=2,3$ and $n\ge5$.
\end{theorem}
\begin{proof}
The complete admissibility classification in [L] leaves exactly
these possibilities and $\GL_2(2)$. For $q=3$ and $n\ge5$, every
nonzero determinant is $1$ or $-1$, so
$\GL_n(3)=H_n(3)$. The complete marking and intersection proof in
[H] therefore give the result in one step, including dimension six.
No intersection property is inferred merely by reducing the integral
generators modulo three.

For $q=2,n\ge5$, we have
$\GL_n(2)=\SL_n(2)=\PSL_n(2)$. The retained simple-group
classification and the intersection criterion give a rank-three
string C-marking. Finally $\GL_2(2)\cong S_3$. Its involutions
are exactly its three transpositions, and no two distinct ones
commute. A rank-three string C-marking requires distinct commuting
end generators, so none exists. All other cases were already
excluded even from admissibility. This proves [T, Theorem 3.2.2].
\end{proof}

The remaining projective polytope assertion is already proved in
[L, Section 5]: $\PGL_n(q)$ has rank-three regular polytopes precisely
for
\[
 n=2,\ q\ne2;\qquad n\ge4,\ \gcd(n,q-1)=2;
 \qquad n\ge5,\ \gcd(n,q-1)=1.
\]
That proof covers the small $S_4$ case explicitly and excludes $S_3$.
Its parameter domains agree with [T, Corollary 3.3.1]. The two
generation classifications in [T, Theorems 3.2.1 and 3.3.1] are
exactly the two classifications proved in [L].

\section{Source qualification and attribution}
The principal statements here belong to Markovskaya and the
collaborators credited in her thesis. The contribution of this
component is their deduction from the shared centre calculation,
central-lift lemma and finite determinant-sign construction. In
particular, the proof of the linear polytope theorem reuses [H]
instead of a new list of intersections in each dimension. This is
not a new classification of simple groups or a claim that a previously
open conjecture has been resolved.

The intended dimension domain is $n\ge2$, as in the source tables.
The literal centre assertion on [T, p.~69] needs the correction
proved in [Z]: even $n$ does not always imply $Z(H)=Z(\SL_n(q))$.
For even $n\ge6$, Corollary 4.2.1's full-centre quotient is
$\PSL_n(q)$ when $e$ is even and is the stated index-two extension
when $e$ is odd. Its polytope conclusion and type $[4p,n]$ hold in
both cases by [H]. For odd $n\ge5$, the two source corollaries
4.2.2 and 4.2.3 follow literally from [Z] and [H]. The defective
four-dimensional involution and the missing square in the
seven-dimensional cancellation are recorded and corrected in those
components. We recover corrected principal results, not every
printed intermediate formula verbatim.

No new computation is needed for this note. Its deductions use the
written companion proofs and the explicit external classifications
listed above. It is not a Lean formalization. Whole-thesis admission
requires the separate complete source ledger and independent review;
neither a compression factor nor a whole-original admission is
granted by this note alone.

\begin{thebibliography}{9}\small
\bibitem{T}[T] I.~A.~Markovskaya, \emph{Linear groups over finite fields
and over the ring of integers as automorphism groups of regular
polytopes of rank 3}, thesis, Siberian Federal University, 2025,
80-page supplied edition; supervisor Ya.~N.~Nuzhin.
\bibitem{I}[I] \emph{Integral linear groups generated by three involutions},
complete reviewed recovery of Chapter 2 of [T].
\bibitem{L}[L] \emph{Three involutions in finite general linear groups},
complete reviewed recovery of Markovskaya--Nuzhin (2025), with its
common paper \emph{Root parameters and fixed-space obstructions}.
\bibitem{Z}[Z] \emph{Central quotients of the determinant-sign subgroup},
complete reviewed N012 component.
\bibitem{H}[H] \emph{Involutions and intersections in groups with
determinant $\pm1$}, complete reviewed N012 component.
\end{thebibliography}
\end{document}
